\subsection{Jeu de Données}
\label{sec:dataset}

\subsubsection{Choix du corpus et justification}

Entraîner un modèle auto-supervisé requiert un corpus où les
invariants sémantiques sont connus et contrôlables. Nous optons pour
une approche de \emph{compilation croisée contrôlée} : plutôt que
d'utiliser des malwares réels (corpus bruités, déséquilibrés, soumis
à des contraintes légales), nous construisons un jeu de données
synthétique à partir de la suite \textbf{GNU Coreutils
v9.10} \cite{coreutils9_10}, un ensemble de $106$ utilitaires POSIX au
comportement sémantique parfaitement documenté.

Ce choix repose sur une propriété fondamentale : deux compilations
différentes d'un même programme \emph{partagent la même sémantique}
malgré des différences syntaxiques profondes (ordonnancement des
blocs, choix des registres, inlining, déroulage de boucles). Le
modèle est ainsi contraint d'apprendre des représentations
\emph{invariantes à la compilation} — exactement la propriété
recherchée pour la détection de variants de malwares.

\subsubsection{Génération du corpus}

Le pipeline de génération est entièrement automatisé via
\texttt{generate\_data.sh}. Les sources sont téléchargées depuis le
dépôt officiel de GNU Coreutils, puis chaque binaire est compilé
selon la grille de variantes définie par le produit cartésien :
\[
  \mathcal{C} \times \mathcal{O}
  = \{\texttt{gcc},\,\texttt{clang}\}
    \times \{\texttt{-O0},\,\texttt{-O1},\,\texttt{-O2},\,\texttt{-O3},\,\texttt{-Os}\}.
\]
Chaque binaire $b \in \mathcal{B}$ (avec $|\mathcal{B}| = 106$) est
donc compilé en $|\mathcal{C}| \times |\mathcal{O}| = 10$ variants
distincts, pour un total de $106 \times 10 = \mathbf{1\,060}$ fichiers
ELF. Deux précautions sont systématiquement appliquées :

\begin{itemize}
  \item \textbf{Stripping complet} (\texttt{--strip-all}) : les
    symboles de débogage, les tables de noms et les sections
    \texttt{.debug\_*} sont supprimés. Cela simule les conditions
    réelles d'analyse de malwares, où les informations symboliques
    sont absentes.
  \item \textbf{Désactivation de la localisation} (\texttt{--disable-nls}) :
    les chaînes de messages localisés sont retirées, réduisant la
    variabilité des données statiques non-sémantiques dans les
    sections \texttt{.rodata}.
\end{itemize}

\subsubsection{Structure du corpus et statistiques}

Le Tableau~\ref{tab:dataset_stats} résume les caractéristiques du
corpus après application du pipeline d'extraction
(Section~\ref{fig:dataset_variants}).

\begin{table}[ht]
\centering
\caption{Statistiques du corpus d'entraînement (GNU Coreutils 9.10,
  après extraction du Bag-of-Paths).}
\label{tab:dataset_stats}
\renewcommand{\arraystretch}{1.3}
\begin{tabular}{lr}
\toprule
\textbf{Caractéristique}                          & \textbf{Valeur}  \\
\midrule
Source                                            & GNU Coreutils 9.10 \\
Nombre de programmes distincts $|\mathcal{B}|$    & 106 \\
Compilateurs $|\mathcal{C}|$                      & 2 (\texttt{gcc}, \texttt{clang}) \\
Niveaux d'optimisation $|\mathcal{O}|$            & 5 (\texttt{-O0} à \texttt{-Os}) \\
Fichiers ELF totaux                               & 1\,060 \\
Symboles conservés                                & Aucun (\texttt{strip --strip-all}) \\
Longueur max de séquence $L$                      & 22000 tokens \\
Taille du vocabulaire $|\mathcal{V}|$             & 371 tokens \\
\bottomrule
\end{tabular}
\end{table}

La Figure~\ref{fig:dataset_variants} illustre la structure de la
grille de compilation et la relation de \emph{co-sémantique} entre
variants d'un même binaire.

% ------------------------------------------------------------------
% Figure : grille de compilation 
% ------------------------------------------------------------------
\begin{figure}[ht]
\centering
\begin{tikzpicture}[
  every node/.style={font=\small},
  elf/.style={
    draw, rounded corners=4pt,
    minimum width=1.55cm, minimum height=0.65cm,
    align=center, font=\scriptsize, line width=0.7pt
  },
  hdrlbl/.style={font=\footnotesize\bfseries},
  arr/.style={-Stealth, line width=0.7pt, gray!60},
]

% ── Pas horizontal entre colonnes
\def\dx{2}

% ── En-têtes colonnes (niveaux d'optimisation)
\foreach \j/\opt in {
  1/\texttt{-O0}, 2/\texttt{-O1}, 3/\texttt{-O2},
  4/\texttt{-O3}, 5/\texttt{-Os}}{
  \node[hdrlbl, gray!70!black] at (\dx*\j, 2.6) {\opt};
}

% ── Source
\node[draw=gray!50, fill=gray!10, rounded corners=4pt,
      minimum width=1.2cm, minimum height=0.9cm,
      align=center, font=\small\ttfamily]
  (src) at (0, 2.6) {ls\\source};
\draw[arr] (src.east) -- ++(0.4, 0);

% ── Ligne gcc (y = 3.5)
\node[hdrlbl, anchor=east, colCtx!80!black] at (1.1, 3.5) {\texttt{gcc}};
\foreach \j/\col in {
  1/colCtx!18, 2/colCtx!28, 3/colCtx!40, 4/colCtx!52, 5/colCtx!28}{
  \pgfmathsetmacro{\optidx}{int(\j - 1)}
  \node[elf, fill=\col, draw=colCtx!65]
    (gcc\j) at (\dx*\j, 3.5)
    {\texttt{ls\_gcc\_O\optidx}};
}

% ── Accolade au-dessus de gcc (bien au-dessus des cellules)
\draw[decorate, decoration={brace, amplitude=6pt},
      colCtx!60!black]
  (gcc1.north west) -- (gcc5.north east)
  node[midway, above=8pt, font=\tiny, colCtx!70!black,
       text width=6cm, align=center]
    {même sémantique, structures syntaxiques différentes};

% ── Ligne clang (y = 1.7)
\node[hdrlbl, anchor=east, colTgt!80!black] at (1.1, 1.7) {\texttt{clang}};
\foreach \j/\col in {
  1/colTgt!18, 2/colTgt!28, 3/colTgt!40, 4/colTgt!52, 5/colTgt!28}{
  \pgfmathsetmacro{\optidx}{int(\j - 1)}
  \node[elf, fill=\col, draw=colTgt!65]
    (clang\j) at (\dx*\j, 1.7)
    {\texttt{ls\_clang\_O\optidx}};
}

% ── Annotations strip (entre les deux lignes)
\node[font=\tiny, gray!65] at (\dx*6, 2.65)
  {\texttt{strip -{}-strip-all}};

% ── Légende bas de figure
\node[draw=colCtx!65, fill=colCtx!15, rounded corners=2pt,
      font=\tiny\bfseries, inner sep=3pt,
      text=colCtx!80!black] (lgcc) at (\dx*2.8, 0.65) {\texttt{gcc}};
\node[draw=colTgt!65, fill=colTgt!15, rounded corners=2pt,
      font=\tiny\bfseries, inner sep=3pt,
      text=colTgt!80!black] (lclang) at (\dx*3.4, 0.65) {\texttt{clang}};
\node[font=\tiny, gray!70!black] at (\dx*3.1, 0.2)
  {10 variants ELF par binaire ($\times\,106 = 1\,060$ au total)};

\end{tikzpicture}
\caption{Grille de compilation pour un binaire représentatif
  (\texttt{ls}). Chaque cellule est un fichier ELF strippé
  indépendant. Les 10 variants d'un même programme partagent
  une sémantique identique mais des représentations binaires
  structurellement distinctes, offrant un signal d'entraînement
  riche pour l'invariance à la compilation.}
\label{fig:dataset_variants}
\end{figure}

\subsubsection{Limitations et perspectives}

Ce corpus présente deux limitations majeures qu'il convient d'anticiper
pour la généralisation à un contexte de détection de malwares réels.

Premièrement, GNU Coreutils est un corpus de \emph{logiciels
bénins}, compilé dans des conditions nominales. Il ne contient ni
obfuscation, ni packing, ni techniques d'anti-analyse. Le modèle
entraîné sur ce corpus apprend des représentations de code
légitime, ce qui constitue une étape préalable à un entraînement
par transfert sur des malwares (\emph{fine-tuning}).

Deuxièmement, l'ensemble des binaires provient d'une unique base
de code (un seul projet, un seul langage source, une seule
architecture cible x86-64). La diversité des patterns sémantiques
est donc plus limitée que dans un corpus multi-famille de malwares.
Une extension naturelle consiste à incorporer d'autres projets
open-source de grande envergure (noyau Linux, OpenSSL, FFmpeg)
afin d'élargir la couverture sémantique du vocabulaire VEX appris.
